\section{Planteamiento del problema}

En los últimos años, el acceso a internet en los hogares del Perú se ha convertido en una necesidad importante para la vida diaria. Este servicio permite realizar actividades educativas, laborales, comerciales y de comunicación. También facilita el acceso a información, trámites digitales y servicios en línea. Por ello, el acceso a internet puede considerarse un indicador relacionado con el desarrollo tecnológico y social del país.

Según los datos del INEI, el porcentaje de hogares con acceso a internet ha aumentado entre los años 2014 y 2024. Sin embargo, este crecimiento no ha sido igual en todo el país. En el área urbana, el acceso a internet es mayor que en el área rural, lo que evidencia una diferencia entre ambos ámbitos de residencia.

Desde el punto de vista matemático, este fenómeno puede estudiarse mediante ecuaciones diferenciales, ya que el porcentaje de hogares con internet cambia con respecto al tiempo. La derivada permite representar la velocidad de crecimiento del acceso a internet durante el periodo estudiado. Además, al tratarse de porcentajes, el crecimiento no puede ser ilimitado, porque el valor máximo posible es 100\%.

Por ello, este trabajo busca modelar el crecimiento del acceso a internet en los hogares peruanos entre 2014 y 2024 mediante ecuaciones diferenciales, usando datos del INEI.
